\section{A triangular trace estimate}\label{sec:trace}

We now compare a cyclic product of three matrices with the Hessian
costs from Section~\ref{sec:matrix}. This is the finite-dimensional
estimate that will control a local triangular integral. Its proof
follows the mixed trace argument in \cite[Section~3]{OpenAI2026};
working over the reals lets us identify the two orientations of the
triangle directly by transposition.

Let \(H_v\), \(v\in\Z/3\Z\), be finite-dimensional real inner product
spaces, and let \(W_v:H_{v+1}\to H_v\). The two edges incident to
vertex \(v\) form the \emph{row star}
\begin{equation}\label{eq:row-star}
 R_v=[\,W_v\quad W_{v-1}^*\,]:
              H_{v+1}\oplus H_{v-1}\longrightarrow H_v,
 \qquad T_v=R_v^*R_v.
\end{equation}
The two Gram matrices \(R_v^*R_v\) and \(R_vR_v^*\) have the same
nonzero eigenvalues with multiplicities, so their \(\Phi\) values
agree.

\begin{proposition}[Mixed trace estimate]\label{prop:trace}
For the maps in \eqref{eq:row-star}, let \(D_v:H_v\to H_v\) be
symmetric with \(\norm{D_v}_{\op}\le1\), and put
\(U_v=R_v^*D_vR_v\). Then \(U_v\) is supported on \(T_v\), and
\begin{equation}\label{eq:trace}
 \abs{\tr(D_0W_0D_1W_1D_2W_2)}
       \le\frac53\sum_{v=0}^2 b_{T_v}(U_v).
\end{equation}
The constant is independent of all three dimensions.
\end{proposition}

\begin{proof}
The identity \(\ker T_v=\ker R_v\) proves the support assertion.
We first place the three quadratic forms on one space, then bound a
cubic trace in an eigenbasis.

\paragraph{A common quadratic form.}
On \(H=H_0\oplus H_1\oplus H_2\), define
\[
 M=\begin{pmatrix}
      0&W_0&W_2^*\\
      W_0^*&0&W_1\\
      W_2&W_1^*&0
    \end{pmatrix},
 \qquad D=\diag(D_0,D_1,D_2).
\]
Both matrices are symmetric and \(\norm D_{\op}\le1\).
Let \(P_v\) be the projection onto \(H_v\), and set
\(A_v=MP_vM\), \(K_v=MP_vDM\).
The only nonzero block row of \(P_vM\) is \(R_v\), with its columns
placed in vertex order. Thus \(A_v\) and \(K_v\) are \(T_v\) and
\(U_v\), respectively, after a permutation of blocks and extension
by zero. In particular \(K_v\) is supported on \(A_v\).
Also \(P_vD=DP_v\), so \(K_v\) is symmetric, and
\[
 0\le A_v\le M^2,\qquad \sum_vK_v=MDM.
\]
Lemma~\ref{lem:order}, followed by the quadratic-form inequality for
a sum of three vectors, gives
\begin{equation}\label{eq:common-form}
 \mathcal E:=\sum_v b_{T_v}(U_v)
 \ge\sum_v b_{M^2}(K_v)
 \ge\frac13b_{M^2}(MDM).
\end{equation}
This comparison is on the relevant supports; \(M\) need not be
invertible.

\paragraph{The cubic trace.}
Choose an orthonormal eigenbasis of \(M\), with eigenvalues \(m_i\)
ordered so that \(q_i=\abs{m_i}\) is nonincreasing. Write \(D_{ij}\)
for the entries of \(D\) in this basis. Since
\((MDM)_{ij}=m_im_jD_{ij}\), \eqref{eq:common-form} gives
\begin{equation}\label{eq:Q}
 \mathcal E\ge\frac12 Q,\qquad
 Q=\sum_{\substack{i,j\\q_i+q_j>0}}
          \frac{q_i^2q_j^2}{q_i+q_j}\abs{D_{ij}}^2.
\end{equation}
Define
\[
 A_i=q_i^3\abs{D_{ii}}^2,\qquad
 E_i=q_i\sum_{j>i}q_j^2\abs{D_{ij}}^2,\qquad
 A=\sum_iA_i,\quad E=\sum_iE_i.
\]
The diagonal terms of \(Q\) sum to \(A/2\). Pairing its off-diagonal
terms and using \(q_i\ge q_j\) for \(i<j\), we obtain
\begin{equation}\label{eq:AE}
 Q=\frac A2+
 \sum_{\substack{i<j\\q_i+q_j>0}}
        \frac{2q_i^2q_j^2}{q_i+q_j}\abs{D_{ij}}^2
 \ge \frac A2+E.
\end{equation}

Expand
\[
 \tr((DM)^3)=\sum_{i,j,k}m_im_jm_kD_{ij}D_{jk}D_{ki}.
\]
Group the triples by their smallest index \(i\) and the number of
times it occurs. If it occurs once, its three possible positions give
\[
 3m_i\sum_{j,k>i}m_jm_kD_{ij}D_{jk}D_{ki}
       =3m_i x^*D_{>i}x,\qquad x_j=m_jD_{ji}.
\]
Here \(j=k\) is allowed, and \(D_{>i}\) is the compression of \(D\)
to indices greater than \(i\). This compression is a contraction,
so the absolute value is at most \(3E_i\).

If the smallest index occurs twice, the contribution is
\[
 3m_i^2D_{ii}\sum_{j>i}m_j\abs{D_{ij}}^2.
\]
The row norm bound \(\sum_j\abs{D_{ij}}^2\le1\) and
Cauchy--Schwarz show that its absolute value is at most
\[
 3q_i^2\abs{D_{ii}}
       \left(\sum_{j>i}q_j^2\abs{D_{ij}}^2\right)^{1/2}
       =3\sqrt{A_iE_i}.
\]
When all three indices equal \(i\), the absolute value is at most
\(A_i\), since \(\abs{D_{ii}}\le1\).
These cases partition all triples, including repeated indices.
Summing and applying Cauchy--Schwarz once more,
\begin{equation}\label{eq:cubic}
 \begin{split}
 \abs{\tr((DM)^3)}
 &\le A+3E+3\sqrt{AE}\\
 &\le \tfrac52A+\tfrac92E
 \le 5(A/2+E)\le5Q.
 \end{split}
\end{equation}
Zero eigenvalues contribute zero throughout.

\paragraph{Returning to the triangle.}
In vertex blocks, a nonzero closed three-step product must visit all
three vertices because \(M\) has zero diagonal blocks. The three
starting points in one orientation give
\(\tau=\tr(D_0W_0D_1W_1D_2W_2)\) by cyclicity, and the reverse
orientation gives the same real trace by transposition.
Thus \(\tr((DM)^3)=6\tau\). Equations~\eqref{eq:Q}
and~\eqref{eq:cubic} yield
\[
 6\abs\tau\le5Q\le10\mathcal E,
\]
which is \eqref{eq:trace}.
\end{proof}
